\begin{problem}[Warmup]{(a)}
    If $\varphi : G \to H$ is a group homomorphism and $G$ is an additive group, need $H$ be an additive group?
\end{problem}
\unsolved

\begin{problem}[Warmup]{(b)}
    Is a ring homomorphism necessarily a group homomorphism?
\end{problem}
\unsolved

\begin{problem}[Warmup]{(c)}
    Suppose we have a group homomorphism between the additive groups of two rings. Is this necessarily a ring homomorphism too?
\end{problem}
\unsolved

\begin{problem}[Warmup]{(d)}
    Check that the determinant function preserves multiplication for two interesting matrices of your choice.
\end{problem}
\unsolved

\begin{problem}[Warmup]{(e)}
    Does the determinant function preserve addition? Try it, for two interesting matrices of your choice.
\end{problem}
\unsolved

\begin{problem}[Warmup]{(f)}
    We proved that the determinant function preserves multiplication for any pair of $2 \times 2$ matrices that are units. Show that this proof works for \emph{any} pair of $2 \times 2$ matrices.
\end{problem}
\unsolved

\begin{problem}[Warmup]{(g)}
    Does the observation in Exercise f mean that the determinant function defined from $M_2(\mathbb{R})$ onto $\mathbb{R}$ is a \emph{ring} homomorphism? \emph{Hint:} What does Exercise e say?
\end{problem}
\unsolved

\begin{problem}[Warmup]{(h)}
    How many polynomials are sent to the polynomial 2 by the derivative homomorphism? Can you describe them all efficiently?
\end{problem}
\unsolved

\begin{problem}[Warmup]{(i)}
    Find a one-to-one group homomorphism $\varphi$ from $\mathbb{Z}_2$ to $\mathbb{Z}_4$. What is the subgroup $\varphi(\mathbb{Z}_2)$?
\end{problem}
\unsolved

\begin{problem}[Warmup]{(j)}
    Let $G$ and $H$ be groups with operations $\circ$ and $*$, respectively, and $\varphi : G \to H$ a group homomorphism. Explain why $\varphi(g \circ h \circ k) = \varphi(g) * \varphi(h) * \varphi(k)$.
\end{problem}
\unsolved

\begin{problem}[Warmup]{(k)}
    Generalize Exercise j to a product of $n$ terms. Does this mean that $\varphi(g^n) = {(\varphi(g))}^n$?
\end{problem}
\unsolved

\begin{problem}{1}
    Show that the function $\varphi : \mathbb{Z}[i] \to \mathbb{Z}$ defined by
    $\varphi(a + bi) = a$ is a group homomorphism (although it is not a ring
    homomorphism).
\end{problem}
\unsolved

\begin{problem}{2}
    Can you find a group homomorphism from $\mathbb{Z}_2 \times \mathbb{Z}_2$ onto the
    multiplicative group $\{1, -1, i, -i\}$?
\end{problem}
\unsolved

\begin{problem}{3}
    Find three distinct group homomorphisms from $\mathbb{Z}_2 \times \mathbb{Z}_2$ onto
    $\mathbb{Z}_2$. How many more homomorphisms exist, if we remove the requirement that
    they be onto?
\end{problem}
\unsolved

\begin{problem}{4}
    Let $G$ be an abelian group (written additively). Define
    \[\psi : G \times G \to G\]
    by $\psi(g, h) = g + h$. Make a two-column chart showing what this function does to
    each element of $G = \mathbb{Z}_3$. Prove that $\psi$ is a homomorphism, for any
    abelian group $G$.
\end{problem}
\unsolved

\begin{problem}{5}
    Show by example that the corresponding homomorphism $\psi$ defined in Exercise 4 need
    \emph{not} be a homomorphism, when the group is not abelian.
\end{problem}
\unsolved

\begin{problem}{6}
    Prove Theorem 27.1: That is, let $\varphi : G \to H$ be a homomorphism between the
    groups $G$ and $H$. Prove that $\varphi(G)$ is a group.
\end{problem}
\unsolved

\begin{problem}{7}
    Suppose that $G$ is an abelian group, and $\varphi$ is a group homomorphism whose
    domain is $G$. Prove that $\varphi(G)$ is abelian.
\end{problem}
\unsolved

\begin{problem}{8}
    Suppose that $R$ and $S$ are rings with unity, and $\varphi : R \to S$ is an onto ring
    homomorphism. Consider the function $\psi : U(R) \to S$ defined by
    $\psi(r) = \varphi(r)$. (That is, $\psi$ is just $\varphi$ restricted to the group of
    units $U(R)$.) Prove that $\psi$ is a group homomorphism from $U(R)$ into $U(S)$.
\end{problem}
\unsolved

\begin{problem}{9}
    We know that $\log : \mathbb{R}^+ \to \mathbb{R}$ is a group homomorphism; see
    Example 27.3. Find a subgroup $G$ of $\mathbb{R}^+$, so that the restricted logarithm
    function maps $G$ onto the subgroup $\mathbb{Z}$ of $\mathbb{R}$. What about a
    subgroup of $\mathbb{R}^+$ that maps onto the subgroup $\mathbb{Q}$ of $\mathbb{R}$?
\end{problem}
\unsolved

\begin{problem}{10}
    Consider the set $G$ of all differentiable real-valued functions. Why is this a group
    under addition? Is the derivative function $D$ a group homomorphism from $G$ to $G$?
    \emph{Warning:} That the function preserves the operation is \emph{not} the issue.
\end{problem}
\unsolved

\begin{problem}{11}
    Consider the homomorphism in Example 27.5. Define the analogous homomorphism from
    $D_4 \to \{1, -1\}$, making use of your work in Exercise 22.1.
\end{problem}
\unsolved

\begin{problem}{12}
    In this problem we consider group homomorphisms from $\mathbb{Z}_n$ to itself.
    \begin{enumerate}
        \item[(a)] Find all homomorphisms from $\mathbb{Z}_5$ into $\mathbb{Z}_5$. Now find
        all homomorphisms from $\mathbb{Z}_6$ into $\mathbb{Z}_6$.
        \item[(b)] Suppose $G$ is a cyclic group of order $n$. How many homomorphisms are
        there from $G$ into $G$? Describe them.
    \end{enumerate}
\end{problem}
\unsolved

\begin{problem}{13}
    In this exercise we generalize the ideas encountered in Exercise 12. A homomorphism
    from a group into itself is called an \textbf{endomorphism} of the group. For an
    abelian group $G$, let $\mathrm{End}(G)$ be the set of all its endomorphisms
    $\varphi : G \to G$.
    \begin{enumerate}
        \item[(a)] Let $\varphi, \psi \in \mathrm{End}(G)$. Define the sum
        $\varphi + \psi$ of these two homomorphisms by
        \[(\varphi + \psi)(g) = \varphi(g) + \psi(g).\]
        Prove that $\varphi + \psi$ is an endomorphism of $G$.
        \item[(b)] Prove that $\mathrm{End}(G)$, under the addition defined in part a, is
        an abelian group.
        \item[(c)] Let $\varphi, \psi \in \mathrm{End}(G)$. Define the product
        $\varphi\psi$ as the ordinary functional composition:
        \[\varphi\psi(g) = \varphi(\psi(g)).\]
        Show that $\mathrm{End}(G)$ is a ring with unity when equipped with this
        multiplication. We call $\mathrm{End}(G)$ the \textbf{endomorphism ring} for $G$.
        \item[(d)] Let $n$ be a positive integer. \emph{Beware:} in this problem we will
        consider $\mathbb{Z}_n$ as an abelian group, and as a commutative ring. Prove that
        the ring $\mathrm{End}(\mathbb{Z}_n)$ is isomorphic as a ring to $\mathbb{Z}_n$.
    \end{enumerate}
\end{problem}
\unsolved
